\chapter{Conclusão}

Este trabalho teve como objetivo realizar uma revisão sistemática da literatura sobre o desenvolvimento e a avaliação de sistemas e interfaces destinados a pessoas com baixa alfabetização informática, buscando identificar e sintetizar métricas, métodos de avaliação, diretrizes e elementos de interface relacionados à usabilidade, à acessibilidade e à Interação Humano-Computador. A pesquisa foi conduzida a partir de um processo sistemático de busca, seleção e análise de estudos científicos, considerando publicações recuperadas nas bases Scopus, IEEE Xplore e Engineering Village.

Ao final do processo de seleção, foram incluídos 79 estudos, sendo 21 provenientes da Scopus, 31 da IEEE Xplore e 27 da Engineering Village. A análise desse conjunto permitiu reunir diferentes perspectivas sobre a interação de pessoas com pouca familiaridade com recursos computacionais, incluindo usuários novatos, pessoas com baixa alfabetização digital e usuários idosos. Embora esses grupos apresentem características distintas, os estudos selecionados oferecem contribuições relevantes para compreender dificuldades e necessidades que podem surgir durante a utilização de sistemas digitais.

Em relação à primeira questão de pesquisa (RQ1), que investigou quais métricas de usabilidade e acessibilidade são utilizadas na avaliação de sistemas e interfaces destinados a pessoas com baixa alfabetização informática, 43 estudos apresentaram respostas diretas, enquanto 29 trouxeram informações relacionadas e 7 não abordaram o assunto. Entre as contribuições identificadas, foram encontradas medidas relacionadas à eficácia, à eficiência e à satisfação dos usuários, como tempo de execução de tarefas, quantidade de erros, taxa de sucesso e avaliações por questionários. Também foram identificados instrumentos e critérios específicos, como o System Usability Scale (SUS), além de avaliações relacionadas à legibilidade, à complexidade visual e à acessibilidade. Os resultados indicam que a avaliação da interação pode envolver diferentes dimensões, não se limitando ao desempenho na execução de tarefas, mas considerando também a compreensão e a experiência dos usuários.

Quanto à segunda questão de pesquisa (RQ2), relacionada aos elementos de interface, diretrizes e princípios de design apresentados na literatura, 63 estudos responderam diretamente e 16 apresentaram informações relacionadas. Nenhum dos estudos selecionados deixou de abordar o tema. As contribuições encontradas apontam para a importância de interfaces com organização visual clara, textos legíveis, linguagem simples, contraste adequado e elementos de navegação consistentes. Também foram identificadas recomendações relacionadas ao tamanho de fontes e botões, à utilização de ícones compreensíveis, à apresentação de instruções, ao fornecimento de feedback e à prevenção e recuperação de erros. Em conjunto, essas recomendações demonstram a relevância de considerar as características e as necessidades dos usuários durante a organização e a apresentação das informações e funcionalidades de um sistema.

A terceira questão de pesquisa (RQ3), que investigou os métodos e as estratégias utilizados para avaliar a usabilidade e a acessibilidade das interfaces, apresentou 67 estudos com respostas diretas e 12 com informações relacionadas. Nenhum dos estudos selecionados deixou de abordar o tema. Entre os métodos identificados estão testes de usabilidade baseados em tarefas, avaliações heurísticas, entrevistas, questionários, observações, protocolos de verbalização durante a execução de atividades e ferramentas automatizadas de verificação de acessibilidade. Também foram encontrados estudos que utilizaram experimentos e análises estatísticas, além de abordagens participativas, nas quais os usuários contribuíram para a identificação de necessidades e para a avaliação das interfaces. Essa variedade evidencia que diferentes métodos podem ser empregados de acordo com os objetivos da avaliação e as características do público investigado.

Por sua vez, a quarta questão de pesquisa (RQ4), voltada às características dos usuários e aos contextos de utilização considerados nos estudos, também apresentou 57 respostas diretas, enquanto 21 trabalhos trouxeram informações relacionadas e 1 não abordou o assunto. A análise identificou pesquisas envolvendo pessoas idosas, usuários novatos, indivíduos com pouca experiência em tecnologias digitais e pessoas com diferentes tipos de deficiência. Entre as características consideradas estão a experiência prévia com computadores e Internet, a escolaridade, a idade e as limitações visuais, auditivas, motoras e cognitivas. Além disso, alguns estudos consideraram aspectos relacionados ao contexto cultural, às condições de saúde e às atividades realizadas pelos participantes. Esses resultados reforçam a importância de compreender o perfil dos usuários e o ambiente em que os sistemas são utilizados, uma vez que diferentes características podem influenciar a forma como ocorre a interação.

De maneira geral, os resultados das quatro questões de pesquisa demonstram que a literatura apresenta contribuições variadas para a compreensão e a avaliação de interfaces destinadas a pessoas com diferentes níveis de familiaridade com recursos digitais. As recomendações de design, as métricas de avaliação e os métodos utilizados aparecem associados a diferentes públicos e contextos de aplicação. Dessa forma, a análise conjunta dos estudos permite observar pontos recorrentes, como a preocupação com a clareza das informações, a facilidade de navegação, a prevenção de erros e a consideração das características dos usuários durante o desenvolvimento e a avaliação de sistemas.

A avaliação da qualidade dos estudos também contribuiu para a análise do material selecionado. Considerando os seis critérios estabelecidos nesta pesquisa, 72 estudos foram classificados como de alta qualidade, correspondendo a 91,1\% do conjunto, enquanto 7 foram classificados como de qualidade média, representando 8,9\%. Não foram identificados estudos classificados como de baixa qualidade. Essa avaliação permitiu caracterizar os trabalhos incluídos e considerar suas diferenças metodológicas durante a interpretação dos resultados. As classificações individuais e as pontuações atribuídas a cada estudo estão apresentadas no apêndice deste trabalho, complementando as informações discutidas ao longo da análise.

Assim, considera-se que o objetivo geral da pesquisa foi alcançado, uma vez que foi possível identificar, analisar e sintetizar contribuições científicas relacionadas às métricas, aos elementos de interface, aos métodos de avaliação e às características dos usuários considerados na literatura. A revisão não teve como finalidade estabelecer uma solução única ou desenvolver uma proposta prática de interface, mas organizar o conhecimento existente e apresentar um panorama das abordagens encontradas nos estudos selecionados.

\section{Limitações da Pesquisa}

Apesar dos resultados obtidos, algumas limitações devem ser consideradas na interpretação desta revisão sistemática da literatura. A primeira está relacionada à seleção das bases de dados utilizadas. Embora Scopus, IEEE Xplore e Engineering Village disponibilizem um conjunto relevante de publicações nas áreas de tecnologia, computação e engenharia, a utilização de outras bases poderia ampliar a quantidade e a diversidade dos estudos identificados.

Outra limitação está relacionada à estratégia de busca e aos critérios de inclusão e exclusão estabelecidos. A seleção dos termos utilizados nas expressões de busca e a definição dos critérios de elegibilidade influenciam diretamente os resultados recuperados. Dessa forma, é possível que estudos relevantes não tenham sido identificados por utilizarem terminologias diferentes para descrever usuários novatos, baixa alfabetização informática, usabilidade ou acessibilidade.

Também deve ser considerada a diversidade dos públicos e contextos presentes nos estudos selecionados. Parte das publicações investiga especificamente pessoas com baixa alfabetização digital ou usuários novatos, enquanto outras abordam principalmente pessoas idosas ou indivíduos com deficiências. Embora esses trabalhos apresentem contribuições relacionadas ao tema investigado, suas características não são necessariamente equivalentes. Portanto, os resultados devem ser interpretados considerando as diferenças entre os grupos estudados, evitando generalizações para todos os usuários com baixa alfabetização informática.

Por fim, a avaliação da qualidade e a classificação dos estudos de acordo com as questões de pesquisa foram realizadas individualmente pelo autor deste trabalho. Apesar da definição prévia de critérios e da aplicação de parâmetros uniformes durante a análise, não houve a participação de um segundo avaliador independente. Essa condição impossibilitou a comparação entre avaliações e a mensuração da concordância entre pesquisadores, podendo permitir a influência de interpretações individuais durante o processo de classificação.

\section{Trabalhos Futuros}

A partir das contribuições identificadas e das limitações apresentadas, algumas possibilidades podem ser consideradas para pesquisas futuras. Uma delas consiste na ampliação da revisão sistemática por meio da inclusão de outras bases de dados, da utilização de novas combinações de termos de busca e da investigação de publicações em diferentes idiomas. Essas iniciativas poderiam ampliar a abrangência do levantamento e permitir a identificação de contribuições que não foram contempladas nesta pesquisa.

Outra possibilidade consiste na realização de estudos empíricos com pessoas que apresentam diferentes níveis de familiaridade com computadores e tecnologias digitais. Pesquisas dessa natureza poderiam investigar como as recomendações identificadas na literatura se relacionam com as dificuldades efetivamente enfrentadas pelos usuários durante a realização de tarefas em sistemas digitais. Também seria possível comparar as experiências de diferentes grupos, considerando aspectos como idade, escolaridade, experiência prévia e necessidades específicas de acessibilidade.

Além disso, pesquisas futuras podem aprofundar a investigação sobre a aplicação conjunta de métricas de usabilidade e acessibilidade, analisando como diferentes instrumentos podem contribuir para uma avaliação mais abrangente da interação. Também podem ser exploradas estratégias de avaliação que considerem não apenas o desempenho na execução de tarefas, mas aspectos relacionados à compreensão das informações, à aprendizagem, à confiança e à autonomia dos usuários.

Por fim, estudos futuros podem investigar de maneira mais aprofundada a influência das características individuais e dos contextos de uso sobre a interação com sistemas digitais. A consideração de diferentes perfis de usuários, incluindo pessoas com baixa alfabetização informática, usuários novatos, idosos e pessoas com deficiência, pode contribuir para ampliar a compreensão das necessidades envolvidas na utilização dessas tecnologias.

Dessa forma, espera-se que a síntese apresentada nesta revisão possa servir como referência para novas investigações sobre usabilidade, acessibilidade e Interação Humano-Computador, contribuindo para o avanço das pesquisas relacionadas à inclusão digital e à interação de pessoas com diferentes níveis de familiaridade com recursos computacionais.